% 问题分析
\section{问题分析}
\label{sec:analysis}

四个问题按"物理核算基准 $\to$ 调度优化 $\to$ 通信联动 $\to$ 资源配置"递进，
且共享同一条可核验数据链：飞行时间、能耗与通信链路均按题给统一规则逐项计算，
任一结果都可由随附脚本复算。

\subsection{问题一：单点往返最大安全载荷与货箱组批}

问题一的核心是为每类运输无人机到各服务区的单点往返任务确定最大安全载荷，并
据此把货箱分组成架次。由等效航程法可知，载荷越大则等效航程越短、单位里程能耗
越高，因此最大安全载荷受质量上限、体积上限与能量上限三者中的最紧者约束，且
远区受地形高差影响爬升能耗大，往往能量上限先达界。求解思路：对三机型与十五
服务区组成的四十五个组合逐一二分求最大安全载荷，再按质量、体积与时限把各服务
区货箱组批为架次，形成后续调度问题"一架次一载荷"的输入。

\subsection{问题二：异构无人机多点多架次调度}

问题二的核心是在 8 架机、14 组共享电池、80 个带时限货箱下求解多目标调度
（架次、完工、能耗、迟到），其中及时性优先。建模思路：以架次为决策单元，先按
时限贪心组批生成初始解并拆出紧急架次，再并行比较多类元启发式，随后以图注意力
强化学习与宽带束搜索精炼 Pareto 前沿。求解思路：调度层按截止期优先，耦合无人机
与电池双资源就绪队列求可行派单；评估执行严格口径，即投递时刻所属服务区必须
与其箱号归属区一致，凡"捎带"结算一律判不可行。

\subsection{问题三：通信约束下运输与中继联合调度}

问题三的核心是无人机全程通信不断链。由自由空间传播加地形遮挡的电波模型，沿
航线以 30\,m 步长逐点判定直连可行性，不可直连点构成需中继集合；中继布设点与
服务时段必须与运输调度联动。求解思路：先在候选位置中确定三个布设点，使初始
布设即覆盖全部需中继采样点且运输无需平移，再由两架中继按时间分片排班，逐项
核算联合完工、中继能耗与机器级排班可行性。

\subsection{问题四：任务分区与资源配置}

问题四的核心是把分区与资源配置作为联合决策。分区应使各组地理与通信需求自洽，
即一组尽量对应一个中继布设点及其覆盖的服务区；资源配置应在满足各组独立作业
时限与容量的前提下核算冗余与库存。求解思路：按中继覆盖结构给出两种自然分区，
对每组以最小机队搜索确定机型与电池数量并叠加冗余，最后对比"按组完全独立配置"
与"任务分区、资源集中池化"两种组织方式的资源缺口，给出库存约束下的工程结论。
